\justifying
\PUAbstractParagraph{This thesis evaluates the seismic performance and practical applicability limits of NBC 205 Ready-to-Use Detailing (RUD) provisions for low-rise reinforced concrete buildings in Nepal. A representative Draft Final NBC 205:2012 building is used as a historical design-check reference. The performance-based comparison considers the NBC 205:2024 Regular model and six plan-irregular configurations (L, L1, L2, T, U, and Non-Parallel), all with three full storeys plus a stair cover. A separate four-storey-plus-stair-cover Regular model is treated as a height-extension case outside the seven-model performance-point comparison.}
\PUAbstractParagraph{The assessment considers modal periods and translational modal mass participation, NBC 105:2025 storey shear, inter-storey drift, centre-of-mass/centre-of-rigidity separation, pushover capacity, performance points, final-step hinge observations, and concrete design-check observations. The Draft Final NBC 205:2012 building is evaluated from the available concrete design-check results, while the nonlinear comparison is based on the NBC 205:2024 analytical results.}
\PUAbstractParagraph{For Soil Type C and $Z=0.40$, the seven-model fundamental periods range from 0.649 to 0.707~s. Mode~1 contributes approximately 79.09--80.32\% of the translational mass in the X direction, while Mode~2 contributes approximately 74.82--80.42\% in the Y direction; cumulative translational participation exceeds 90\% by Mode~4 in X and Mode~5 in Y for every configuration. For the seven-model NBC 205:2024 comparison, the Regular building reaches approximately 532.489 kN at Story 1 under Soil Type C and $Z=0.40$. The largest ULS drift ratio in that demand case is 0.004954 for the U-shaped building, below the 0.025 ULS drift limit. The largest first-peak pushover base shear is approximately 1882.070 kN for the U-shaped building in the Y direction. The largest performance-point base shear is approximately 456.253 kN for the Regular building in the X direction, whereas the largest performance-point roof displacement is approximately 27.072 mm for the U-shaped building in the X direction. GII does not show a clear monotonic relationship with performance-point response; therefore, GII alone is insufficient to characterize seismic performance. The results support interpretation of the investigated configurations but do not justify unrestricted application of RUD provisions to all irregular buildings without project-specific seismic verification.}
